\section{Introducción y formulación del problema}
\label{sec:intro}

La retención de clientes es una de las principalesprioridades estratégicas en las 
empresas de telecomunica-ciones. En mercados altamente competitivos, la pérdida de 
clientes no solo reduce los ingresos mensuales, sino queademás incrementa los 
costos asociados a la adquisiciónde nuevos usuarios y puede deteriorar la 
estabilidad delportafolio de clientes. En este contexto, comprender qué factores 
explican la salida de clientes resulta esencial para diseñar políticas de retención, 
mejorar la experiencia del usuario y optimizar los servicios ofrecidos.

El problema de \emph{churn} se define como la decisión de uncliente de abandonar un servicio o 
cambiar de proveedor. En la práctica, este fenómeno resulta observable a través de 
datos históricos de comportamiento, uso del servicio, duración del contrato y 
variables de facturación. La capacidad de identificar tempranamente a clientes 
con alto riesgo de churn permite a la empresa intervenir de forma proactiva antes de 
que el cliente abandone por completo.

Este proyecto se centra en el análisis del conjunto de datos Telco Customer Churn~\cite{dataset}, 
una base ampliamente usadaen el contexto de la predicción de retención. IBM 
lo presenta como información de ejemplo de una empresa ficticia~\cite{ibm}.

\subsection{Objetivo}
El objetivo principal es comprender el problema de churn desde una perspectiva estadística y 
analítica, formular la tarea de Machine Learning y preparar la base para el 
modelado supervisado en etapas posteriores. 
Asimismo, se busca comprender el contexto del problema y su relevancia empresarial; 
caracterizar el dataset y sus variables relevantes; detectar patrones de churn en los datos;
identificar riesgo de leakage, faltantes y anomalías; así como preparar una base metodológica
sólida para la etapa de modelado predictivo.

\subsection{Preguntas de investigación}
\begin{enumerate}
    \item ¿Qué características presentan diferencias en la proporción de abandono observado?
    \item ¿Las diferencias por contrato se mantienen al comparar grupos de antigüedad similares?
    \item ¿Cómo se relacionan el tipo de internet y los servicios adicionales con el abandono, y qué 
    diferencias permanecen dentro de cada contrato?
    \item ¿Qué problemas de calidad, representatividad y fuga de información deben considerarse antes 
    del modelado?
\end{enumerate}


\subsection{Formulación del problema}
El problema de churn puede representarse como una tarea de clasificación supervisada binaria. 
Dado un conjunto de observaciones $\mathbf{x}_i$ que describen a cada cliente, se 
busca aprender una función $f(\mathbf{x}_i)$ tal que:

\begin{equation}
\hat{y}_i = f(\mathbf{x}_i) \in \{\text{No}, \text{Yes}\}
\end{equation}

donde $\hat{y}_i$ representa la predicción de si el cliente abandona el servicio (\var{Yes}) o no (\var{No}). 

El conjunto de variables de entrada incluye características demográficas, 
condiciones del contrato, servicios contratados, pagos y métricas de uso. La 
variable objetivo, Churn, es binaria y resume la decisión final del cliente respecto 
a su continuidad con el proveedor.

Desde la perspectiva del aprendizaje automático, el reto principal no solo reside en 
la predicción, sino en la correcta identificación de clientes en riesgo. Esto exige 
un análisis robusto de las relaciones entre variables, así como una evaluación 
orientada a métricas sensibles al desbalance de clases.

